\documentclass[11pt]{article}
\usepackage[margin=1in]{geometry}
\usepackage{amsmath,amssymb,amsthm,mathtools}
\usepackage[hidelinks]{hyperref}

\newtheorem{theorem}{Theorem}[section]
\newtheorem{lemma}[theorem]{Lemma}
\newtheorem{proposition}[theorem]{Proposition}
\newcommand{\spE}{s^+}
\newcommand{\sig}{\sigma}

\title{Positive Square Energy of Bridge-Separated Tetracyclic Cacti}
\author{}
\date{25 July 2026}

\begin{document}
\maketitle

\begin{abstract}
Let $G$ be a connected tetracyclic cactus of order $n$.  We prove
$s^+(G)>n$ whenever its four cyclic blocks do not all belong to one connected
shared-cut cluster.  The sharp block-additive doubly nonnegative bound reduces
the problem to cycle multisets $\{3,3,3,q\}$ and
$\{3,3,5,5\}$.  We settle every bridge-separated incidence by adaptive
induced packetization.  Two apparent obstructions disappear only after one
allows a hostile cycle to be sacrificed into path fragments: a shared
$C_3-C_5-C_3$ cluster plus a remote pentagon, and a shared three-triangle
cluster plus a remote $1\pmod4$ cycle.  Arbitrary connector trees and arbitrary
trees attached at arbitrary vertices are allowed.  The case in which all four
cycles form one shared-cut cluster is not asserted here.
\end{abstract}

\section{Preliminaries}

All graphs are finite, simple, and undirected.  If the adjacency eigenvalues
are $\lambda_1,\ldots,\lambda_n$, put
\[
 \spE(G)=\sum_{\lambda_j>0}\lambda_j^2,
 \qquad \sig(G)=\spE(G)-|V(G)|.
\]
We use induced-subgraph superadditivity~\cite{AKMP}: for every vertex partition
$V(G)=V_1\dot\cup\cdots\dot\cup V_k$,
\begin{equation}\label{eq:super}
 \spE(G)\geq\sum_{j=1}^k\spE(G[V_j]).
\end{equation}

Join two cyclic blocks in the \emph{shared-cut graph} when they share a cut
vertex, and call its connected components shared-cut clusters.  In the
block-cut tree, contract the incidence subtree belonging to each shared-cut
cluster to one marked node; every cluster, including every singleton cyclic
block, is marked.  Take the minimal subtree spanning these marked nodes and
suppress only unmarked degree-two nodes.  No marked node is ever suppressed.
The resulting tree is the reduced cluster tree.  Consequently the interior of
every reduced edge contains only bridge blocks and unmarked cut nodes, never an
unassigned cyclic block.

\begin{lemma}[Connector territories]\label{lem:territories}
Let $S_1,\ldots,S_k$ be pairwise vertex-disjoint connected subtrees of the
reduced cluster tree such that every marked cluster node belongs to exactly one
$S_i$.  Then there is a vertex partition into connected induced territories
such that the territory corresponding to $S_i$ retains exactly the cyclic
blocks in the clusters marked in $S_i$.
\end{lemma}

\begin{proof}
Every reduced edge represents a path whose block nodes are all bridge blocks in
the block-cut tree; this follows from the preceding convention that every
shared-cut cluster, including a singleton cycle, is a nonsuppressed marked
node.
Assign the internal vertices of each selected connector path to its selected
subtree.  Between two selected subtrees, cut one actual bridge on their joining
path and assign the two remnants to their respective sides.  At an unmarked
Steiner node, assign the node and its initial path segments to one selected
subtree and cut one bridge on every other incident segment.  Every component
outside the minimal cyclic hull meets it at a unique vertex, since two
attachments would create a cycle in the block-cut tree; assign that entire
hanging tree to the territory containing its attachment.  The resulting sets
are connected.  Their induced subgraphs retain precisely the assigned cycles.
\end{proof}

We will also split a cyclic block by assigning a proper path of it to each of
two territories.  Crossing edges are harmless in \eqref{eq:super}.

We quote the following previously proved packet estimates
\cite{Bicyclic,Tricyclic}.  Each permits
arbitrary tree attachments; the mixed bound permits arbitrary incidence and
connectors, while \eqref{eq:pack4} and \eqref{eq:pack5} have their displayed
shared-cut hypotheses.  Write $T=C_3$,
$P=C_5$, and, for $q\equiv1\pmod4$,
\[
 \delta_q=\sec(\pi/q)-1<1,
 \qquad \delta_5=\sqrt5-2.
\]
Then
\begin{align}
 \sig(T)&>0, & \sig(C_q)&\geq-\delta_q,\label{eq:pack1}\\
 \sig(TT)&>1, & \sig(TC_q)&>1-\delta_q,\label{eq:pack2}\\
 \sig(H)&\geq0 &&\text{for every bicyclic or tricyclic cactus }H,\label{eq:pack3}\\
 \sig(TTC_q)&>2-\delta_q
 &&\text{if the two triangles share a cut vertex},\label{eq:pack4}\\
 \sig(TPP)&>6-2\sqrt5
 &&\text{if the three cycles form one shared-cut cluster}.\label{eq:pack5}
\end{align}
If every cycle is $3\pmod4$ and the cycle-packing number is at most two,
the Sachs half-plane theorem~\cite{PackingTwo} gives $D=s^+-s^->0$.  Thus a connected
$r$-cyclic such cactus has
\begin{equation}\label{eq:favorable}
 \sig(G)>r-1.
\end{equation}
These results are proved in the companion manuscripts on all bicyclic and all
tricyclic cacti.

\section{The exact residual reduction}

For a cactus with bridge-block count $b$ and cycle lengths
$\ell_1,\ldots,\ell_4$, block counting gives
\[
 b+\sum_{j=1}^4\ell_j=n+3.
\]
The sharp cactus DNN constant~\cite{DNN,LTZ} implies
\begin{equation}\label{eq:dnn}
 \spE(G)-n\geq3-\sum_{j=1}^4\epsilon_{\ell_j},
\end{equation}
where $\epsilon_\ell=0$ for even $\ell$ and
$\epsilon_\ell=\ell\tan^2(\pi/(2\ell))$ for odd $\ell$.  Moreover
\[
 \epsilon_3=1,\qquad \epsilon_5=5-2\sqrt5,
 \qquad \epsilon_5+\epsilon_7<1,
\]
and $\epsilon_\ell$ strictly decreases through odd lengths.  Hence
\eqref{eq:dnn} proves the theorem unless the cycle multiset is
\begin{equation}\label{eq:residual}
 \{3,3,3,q\}\quad(q\geq3),
 \qquad\text{or}\qquad \{3,3,5,5\}.
\end{equation}

\section{Three triangles and one further cycle}

\begin{proposition}\label{prop:333q}
Let the cyclic blocks of $G$ have lengths $3,3,3,q$, where $q\geq3$.  If
the shared-cut graph is disconnected, then $\spE(G)>n$.
\end{proposition}

\begin{proof}
If $q$ is even, take the shared-cut clusters as induced territories using
Lemma~\ref{lem:territories}.  The territory containing the even cycle is at
most tricyclic and has nonnegative surplus by \eqref{eq:pack3} (and is
bipartite with zero surplus if it contains no triangle).  Since the shared-cut
graph is disconnected, another territory contains one, two, or three
triangles.  Its cycle-packing number is at most two, so its surplus is strictly
positive by \eqref{eq:favorable}.  Every remaining territory has nonnegative
surplus.  Equation~\eqref{eq:super} proves the claim.

If $q\equiv3\pmod4$, partition along the bridge connectors between shared-cut
clusters.  Every territory contains only favorable cycles and has cycle
packing at most two within each nontrivial cluster; singleton territories are
favorable unicyclic cacti.  Equations \eqref{eq:favorable} and
\eqref{eq:super} give the result.

Suppose $q\equiv1\pmod4$, and denote its block by $Q$.  Up to permutation of
the triangles, the shared-cut component partitions are
\[
 T|T|T|Q,\quad TT|T|Q,\quad TQ|T|T,
 \quad TTT|Q,\quad TT|TQ,\quad TTQ|T.
\]
We treat them in turn.

For $T|T|T|Q$, choose a triangle whose marked node has minimum distance from
$Q$ in the reduced cluster tree.  The unique $Q$--$T$ path then contains no
other marked node.  Use this path as one connected subtree and the other two
triangle nodes as singleton subtrees; Lemma~\ref{lem:territories} realizes the
territories even when the path passes through an unmarked Steiner node.  By
\eqref{eq:pack1}--\eqref{eq:pack2}, the total surplus is
strictly greater than $1-\delta_q>0$.  The same direct accounting settles
$TQ|T|T$.  The partition $TT|TQ$ has surplus greater than
$1+(1-\delta_q)>0$.

For $TT|T|Q$, call the three marked nodes $A=TT$, $B=T$, and $C=Q$.  If the
$B$--$C$ path avoids $A$, Lemma~\ref{lem:territories} gives a $TQ$ packet and
a separate $TT$ packet.  Otherwise $A$ lies on that path; the connected
$A$--$C$ subtree gives a $TTQ$ packet separated from $B$.  The resulting
$TTQ$ territory retains the shared pair of triangles, so
\eqref{eq:pack4}, together with the remaining triangular packet, gives total
surplus greater than $2-\delta_q>0$.

For $TTQ|T$, the tricyclic territory has nonnegative surplus by
\eqref{eq:pack3}, and the singleton triangle has strictly positive surplus.

It remains to treat $TTT|Q$.  Let $w\in V(Q)$ be the vertex through which the
block-tree route from $Q$ to the triangular cluster leaves $Q$.  Choose
$v\in V(Q)\setminus\{w\}$.  Let $F$ consist of $v$ together with every
component of $G-E(Q)$ attached to $Q$ only at $v$, and put all other vertices
into $H$.  The graph $Q-v$ is a path containing $w$, so $H$ is
connected; its only cycles are the three triangles in one shared-cut cluster.
Thus its cycle-packing number is at most two and
\[
 \spE(H)>|H|+2
\]
by \eqref{eq:favorable}.  Since $F$ is a tree,
$\spE(F)=|F|-1$.  Therefore
\[
 \spE(G)\geq\spE(H)+\spE(F)>|H|+2+|F|-1=n+1.
\]
All component partitions have now been covered.
\end{proof}

\section{Two triangles and two pentagons}

The only delicate packetization is recorded separately.

\begin{lemma}[Middle-pentagon sacrifice]\label{lem:sacrifice}
Suppose one shared-cut cluster consists of a pentagon $P$ and two disjoint
triangles $T_1,T_2$ meeting $P$ at distinct vertices $x_1,x_2$.  Suppose a
second pentagon $Q$ is bridge-separated from this cluster.  Then the vertices
can be partitioned into connected induced territories which are either
\begin{enumerate}
\item a $TT$ bicyclic cactus and a pentagonal unicyclic cactus, or
\item a triangular unicyclic cactus and a $TP$ bicyclic cactus.
\end{enumerate}
In either case $\spE(G)>n$.
\end{lemma}

\begin{proof}
Let $z$ be the first vertex of the shared cluster met by the unique block-tree
route from $Q$.  Include $Q$, its entire connector, and all branches hanging
from that connector on the $Q$ side.

If $z\in V(P)\setminus\{x_1,x_2\}$, put $z$ with $Q$ and put
$P-z$ together with both triangles on the other side.  The latter is connected
because $P-z$ is a path containing $x_1,x_2$.  The two induced territories
have cycle multisets $\{Q\}$ and $\{T_1,T_2\}$, respectively.  Hence their surplus sum is
strictly greater than $-\delta_5+1=3-\sqrt5>0$.

Otherwise $z$ lies on one triangle, say $T_1$; this includes $z=x_1$ (which
also lies on $P$) and a
connector entering through any branch rooted on $T_1$.  Put all of $T_1$, the
connector, and $Q$ on one side.  Put $P-x_1$ and $T_2$ on
the other side.  The first induced territory is a mixed $TP$ bicyclic cactus;
the second is triangular unicyclic because $P-x_1$ is a path containing
$x_2$.  Their surplus sum is strictly greater than $1-\delta_5>0$.  The case
in which the route enters through $T_2$ is symmetric.

Every remaining bridge-tree component meets this retained skeleton at only one
vertex; otherwise the block-cut graph would contain a cycle.  Assign it wholly
to the territory containing that attachment.  Thus both territories remain
connected and induced.  In all cases the middle pentagon is deliberately
split into proper path fragments, so it creates no uncounted cycle.
\end{proof}

\begin{proposition}\label{prop:3355}
Let the cyclic blocks of $G$ have lengths $3,3,5,5$.  If the shared-cut graph
is disconnected, then $\spE(G)>n$.
\end{proposition}

\begin{proof}
Write again $T=C_3$ and $P=C_5$.  Up to equal-cycle permutation, the possible
shared-cut component partitions are
\[
 TTP|P,\quad TPP|T,\quad TT|PP,\quad TP|TP,
 \quad TT|P|P,\quad PP|T|T,\quad TP|T|P,
 \quad T|T|P|P.
\]

For $TPP|T$, use \eqref{eq:pack5} and the strict triangular surplus.  For
$TT|PP$ and $TP|TP$, the packet sums are respectively greater than $1$ and
$2(1-\delta_5)$.  For the three-cluster partitions, direct cluster territories
give
\begin{align*}
 TT|P|P &: &\sig(G)&>1-2\delta_5=5-2\sqrt5>0,\\
 PP|T|T &: &\sig(G)&>0,\\
 TP|T|P &: &\sig(G)&>1-2\delta_5>0.
\end{align*}

For $T|T|P|P$, consider the reduced tree spanning the four cycle nodes.  If a
triangle is a leaf, cut its first bridge: the triangular packet has positive
surplus and the remaining tricyclic cactus has nonnegative surplus.  If neither
triangle is a leaf, the only leaves are the two pentagons.  A finite tree with
exactly two leaves is a path, so the cycle order is $P-T-T-P$.  Cut between the
triangles to obtain two induced $TP$ bicyclic territories, whose total surplus
is greater than $2(1-\delta_5)>0$.

Finally consider $TTP|P$.  If the two triangles in the triple cluster share a
cut vertex, \eqref{eq:pack4} gives total surplus greater than
$2-2\delta_5>0$.  Otherwise the shared-cut incidence tree is necessarily
$T-P-T$, with the pentagon in the middle and distinct triangle cuts.  This is
exactly Lemma~\ref{lem:sacrifice}.  The enumeration is exhaustive.
\end{proof}

\section{Main theorem}

\begin{theorem}\label{thm:main}
Let $G$ be a connected tetracyclic cactus on $n$ vertices.  If its four cyclic
blocks do not all belong to one connected shared-cut cluster, then
\[
 \boxed{\spE(G)>n}.
\]
\end{theorem}

\begin{proof}
The DNN estimate \eqref{eq:dnn} settles every cycle multiset except
\eqref{eq:residual}.  Proposition~\ref{prop:333q} settles the first residual
family and Proposition~\ref{prop:3355} settles the second.  Every construction
uses vertex-disjoint induced territories, so arbitrary crossing bridge edges
are discarded only through the valid superadditivity inequality
\eqref{eq:super}; no edge-monotonicity is used.
\end{proof}

\section{Scope}

The theorem does not cover a tetracyclic cactus in which all four cycles form
one connected shared-cut cluster.  In particular, it does not claim the full
tetracyclic-cactus case.  An exact repository census finds twenty
quotient-isomorphism types for the shared $\{3,3,5,5\}$ residual, and the
direct coefficientwise two-pentagon phase certificate succeeds for only one;
see \texttt{research/c3-c3-c5-c5-shared-cluster-census-2026-07-25.md}.
The other cases remain unresolved by that certificate.

\begin{thebibliography}{9}
\bibitem{AKMP}
S.~Akbari, H.~Kumar, B.~Mohar, and S.~Pragada,
\emph{A linear lower bound for the square energy of graphs},
Electron. J. Combin. 32(3) (2025), P3.53.
\bibitem{AKMPZ}
S.~Akbari, H.~Kumar, B.~Mohar, S.~Pragada, and S.~Zhang,
\emph{Refinement of a conjecture on positive square energy of graphs},
arXiv:2506.07264, 2025.
\bibitem{Bicyclic}
Companion manuscript, \emph{Positive Square Energy of Every Bicyclic Cactus},
repository path \texttt{all-bicyclic-cacti/paper.tex}.
\bibitem{Tricyclic}
Companion manuscript, \emph{Positive Square Energy of Every Tricyclic Cactus},
repository path \texttt{all-tricyclic-cacti/paper.tex}.
\bibitem{DNN}
Companion manuscript, \emph{The Optimized Liu--Tang--Zhang Constant for Cactus
Graphs}, repository path \texttt{sharp-cactus-dnn/paper.tex}.
\bibitem{PackingTwo}
Companion manuscript, \emph{Square Energy from Cycle Packing and Block-Graph
Territories}, repository path \texttt{packing-two-square-energy/paper.tex}.
\bibitem{LTZ}
Y.~Liu, Q.~Tang, and S.~Zhang,
\emph{The positive and negative square-energy conjecture},
arXiv:2607.18031, 2026.
\end{thebibliography}
\end{document}
